\documentclass[11pt]{article}
\usepackage[a4paper,margin=1in]{geometry}
\usepackage{amsmath,amssymb,amsthm,booktabs}
\usepackage[T1]{fontenc}
\usepackage{lmodern}
\usepackage[hidelinks]{hyperref}
\setlength{\emergencystretch}{3em}

\theoremstyle{plain}
\newtheorem{theorem}{Theorem}
\newtheorem{lemma}[theorem]{Lemma}
\newtheorem{proposition}[theorem]{Proposition}
\newtheorem{corollary}[theorem]{Corollary}
\theoremstyle{definition}
\newtheorem{definition}[theorem]{Definition}
\theoremstyle{remark}
\newtheorem{remark}[theorem]{Remark}

\title{A Proof of Conjecture 00000000117:\
Pandigital Primes in Every Base from 2 to 16}
\author{%
  Feiyu\thanks{Independent researcher.}
  \and
  Claude (Opus 5.5)\thanks{Anthropic.}%
}
\date{2026-10-02}

\begin{document}
\maketitle

\begin{abstract}
Conjecture 00000000117 asserts that for every base $b\in[2,16]$ there is a
$b$-ary pandigital prime, a prime whose base-$b$ expansion contains each
digit $0,\ldots,b-1$ at least once. We exhibit one in every base: the smallest
one, as listed in OEIS A185122. We verify the digits of each witness, and we
prove its primality with Pratt certificates checked against the Lucas
primality criterion. Everything is formalized in Lean~4 against Mathlib.
\end{abstract}

\section{The conjecture as stated}

The original text of conjecture 00000000117 reads:

\begin{quote}
Conjecture: For every base $b \in [2, 16]$ there exists a $b$-ary pandigital
prime (whose expansion contains each of the digits $0, \ldots, b-1$ at least
once).
\end{quote}

\section{Reading of the statement}

For an integer $b\ge2$, the base-$b$ expansion of a positive integer $n$ is
its usual digit string without leading zeros. $n$ is \emph{$b$-ary
pandigital} if every $d\in\{0,\ldots,b-1\}$ occurs in it. The conjecture
asserts that for each integer $b$ with $2\le b\le16$ some prime is $b$-ary
pandigital. The statement is finite and unambiguous.

\section{Result}

\begin{theorem}\label{thm:main}
For each $b\in\{2,\ldots,16\}$ the number $p_b$ of Table~\ref{tab:wit} is a
$b$-ary pandigital prime. Hence Conjecture 00000000117 holds.
\end{theorem}

\begin{table}[ht]
\centering
\begin{tabular}{rrl}
\toprule
$b$ & $p_b$ & base-$b$ expansion of $p_b$\\
\midrule
2 & 2 & \texttt{10}\\
3 & 11 & \texttt{102}\\
4 & 283 & \texttt{10123}\\
5 & 3319 & \texttt{101234}\\
6 & 48761 & \texttt{1013425}\\
7 & 863231 & \texttt{10223465}\\
8 & 17119607 & \texttt{101234567}\\
9 & 393474749 & \texttt{1012346785}\\
10 & 10123457689 & \texttt{10123457689}\\
11 & 290522736467 & \texttt{1022345689A7}\\
12 & 8989787252711 & \texttt{101234568A79B}\\
13 & 304978405943587 & \texttt{10123456789ABC}\\
14 & 11177758345241723 & \texttt{10123456789CDAB}\\
15 & 442074237951168419 & \texttt{10223456789ADBCE}\\
16 & 18528729602926047181 & \texttt{10123456789ABEFCD}\\
\bottomrule
\end{tabular}
\caption{The witnesses: the smallest $b$-ary pandigital primes (OEIS
A185122). Digits $10,\ldots,15$ are written \texttt{A}--\texttt{F}.}
\label{tab:wit}
\end{table}

\section{Proof}

\paragraph{Digits.} Reading off the expansions in Table~\ref{tab:wit}, each
contains every digit $0,\ldots,b-1$. (In Lean, $p_b$ is written as
$\sum_i d_i b^i$ for the listed digit string, every digit is checked to be
$<b$ with a nonzero leading digit, so the string is \emph{the} base-$b$
expansion of $p_b$, and membership of each digit is then checked directly.)

\paragraph{Primality.} The witnesses $p_2,\ldots,p_7$ are below $10^6$ and
their primality is checked by trial division. For the others we use the
classical criterion of Lucas, in the form proved in Mathlib.

\begin{lemma}[Lucas]\label{lem:lucas}
Let $p\ge2$ and $a$ be integers with $a^{p-1}\equiv1\pmod p$ and
$a^{(p-1)/q}\not\equiv1\pmod p$ for every prime $q\mid p-1$. Then $p$ is
prime.
\end{lemma}

\begin{proof}
The order of $a$ in $(\mathbb Z/p\mathbb Z)^\times$ divides $p-1$ but no
$(p-1)/q$, so it equals $p-1$. Hence $(\mathbb Z/p\mathbb Z)^\times$ has
$p-1$ elements and every nonzero residue is a unit, so $p$ is prime.
\end{proof}

To apply Lemma~\ref{lem:lucas} one needs all prime factors of $p-1$, which
in turn are certified prime in the same way unless they are small: this is a
Pratt certificate \cite{pratt}. Table~\ref{tab:cert} lists every certificate
used. For each row, $p-1$ equals the displayed product, the witness $a$
satisfies the two congruences of Lemma~\ref{lem:lucas}, and each factor is
either below $10^6$ (prime by trial division) or itself a row of the table.

\begin{table}[ht]
\centering
\small
\begin{tabular}{rrl}
\toprule
$p$ & $a$ & factorization of $p-1$\\
\midrule
1065901 & 6 & $2^{2} \cdot 3 \cdot 5^{2} \cdot 11 \cdot 17 \cdot 19$ \\
1222829 & 2 & $2^{2} \cdot 347 \cdot 881$ \\
6395407 & 3 & $2 \cdot 3 \cdot 1065901$ \\
17119607 & 5 & $2 \cdot 7 \cdot 1222829$ \\
67222271 & 7 & $2 \cdot 5 \cdot 2087 \cdot 3221$ \\
335476601 & 6 & $2^{3} \cdot 5^{2} \cdot 47 \cdot 89 \cdot 401$ \\
393474749 & 2 & $2^{2} \cdot 691 \cdot 142357$ \\
10123457689 & 13 & $2^{3} \cdot 3^{2} \cdot 4091 \cdot 34369$ \\
290522736467 & 2 & $2 \cdot 433 \cdot 335476601$ \\
8989787252711 & 11 & $2 \cdot 5 \cdot 11 \cdot 131617 \cdot 620933$ \\
35485169204621 & 2 & $2^{2} \cdot 5 \cdot 41 \cdot 197 \cdot 317 \cdot 347 \cdot 1997$ \\
50829734323931 & 2 & $2 \cdot 5 \cdot 1129 \cdot 13873 \cdot 324529$ \\
132708276772139 & 2 & $2 \cdot 19 \cdot 47 \cdot 170179 \cdot 436627$ \\
304978405943587 & 2 & $2 \cdot 3 \cdot 50829734323931$ \\
11177758345241723 & 2 & $2 \cdot 13 \cdot 6395407 \cdot 67222271$ \\
442074237951168419 & 2 & $2 \cdot 6229 \cdot 35485169204621$ \\
18528729602926047181 & 6 & $2^{2} \cdot 3 \cdot 5 \cdot 13 \cdot 179 \cdot 132708276772139$ \\
\bottomrule
\end{tabular}
\caption{Pratt certificates, in increasing order of $p$.}
\label{tab:cert}
\end{table}

The modular powers $a^{e}\bmod p$ with $p$ up to $1.85\cdot10^{19}$ are
computed by binary exponentiation, so each check takes about $65$
multiplications modulo $p$.

\section{Formalization}

The accompanying Lean~4 project (\texttt{lean/Submission/Basic.lean}) states
the conjecture and proves it against Mathlib.

\begin{itemize}
  \item \texttt{IsPandigital b p} --- every $d<b$ occurs in
    \texttt{Nat.digits b p}, Mathlib's base-$b$ expansion.
    \texttt{ConjectureHolds} --- every $b$ with $2\le b\le16$ has a prime
    $p$ with \texttt{IsPandigital b p}.
  \item \texttt{powMod} --- binary exponentiation modulo $m$, with
    \texttt{powMod\_eq}: it equals $a^e\bmod m$ whenever $e<2^{\mathrm{fuel}}$.
  \item \texttt{mem\_of\_prime\_dvd\_prod} --- a prime dividing
    $\prod q_i^{e_i}$ (with all $q_i$ prime) is one of the $q_i$.
  \item \texttt{prime\_of\_cert} --- Lemma~\ref{lem:lucas} with a factored
    $p-1$, derived from Mathlib's \texttt{lucas\_primality}.
  \item \texttt{prime\_1065901}, \ldots, \texttt{prime\_18528729602926047181}
    --- the rows of Table~\ref{tab:cert}.
  \item \texttt{digits\_b}, \texttt{isPandigital\_b} for $b=2,\ldots,16$ ---
    the expansions of Table~\ref{tab:wit}, via Mathlib's
    \texttt{Nat.digits\_ofDigits}.
  \item \texttt{conjecture\_00000000117} --- \texttt{ConjectureHolds}.
\end{itemize}

All numerical checks are carried out by the Lean kernel, with
\texttt{decide +kernel} and \texttt{norm\_num}. The kernel evaluates numeral
arithmetic natively, and \texttt{powMod} keeps every intermediate value below
$p^2$. The toolchain is \texttt{leanprover/lean4:v4.33.1}, Mathlib is pinned to
\texttt{v4.33.1}, and \texttt{lake build} succeeds in a few seconds. The
project contains no \texttt{sorry}, no \texttt{admit} and no
\texttt{native\_decide}, and introduces no axioms:
\begin{center}
\texttt{\#print axioms conjecture\_00000000117}
\end{center}
reports exactly \texttt{propext}, \texttt{Classical.choice} and
\texttt{Quot.sound}.

\section{Remarks}

\begin{remark}[why the witnesses repeat a digit]
For $b\ge4$, no pandigital prime has exactly $b$ digits. If each digit occurs
once, then $n\equiv0+1+\cdots+(b-1)=b(b-1)/2\pmod{b-1}$. For even $b$ this
is $\equiv0$, so $b-1\mid n$. For odd $b$ it is $\equiv(b-1)/2$, so
$(b-1)/2\mid n$. In both cases $n$ has a proper divisor greater than $1$
once $b\ge4$. This is why every witness in Table~\ref{tab:wit} with $b\ge4$
has $b+1$ digits.
\end{remark}

\begin{remark}
The witnesses are the smallest pandigital primes, as recorded in OEIS A185122 \cite{oeis};
a computer search confirms this, but minimality is not needed and is not part
of the Lean development. Any pandigital prime in each base would do.
\end{remark}

\begin{thebibliography}{9}
\bibitem{pratt} V.~R.~Pratt, \emph{Every prime has a succinct certificate},
  SIAM J. Comput. 4 (1975), 214--220.
\bibitem{oeis} OEIS Foundation, \emph{The On-Line Encyclopedia of Integer
  Sequences}, sequence A185122 (minimum pandigital prime in base $n$).
\end{thebibliography}

\end{document}
